\section{Related Work}
\subsection{Complex Instruction Generation}

% Instruction tuning is essential for aligning Large Language Models (LLMs) with user intentions. Initial endeavors in this field primarily focused on the curation and refinement of existing open-source NLP datasets \cite{wang2023far,ding2023enhancing}. With the importance of instruction quality recognized, manual annotation methods emerged~\cite{wang2023far,zhou2024lima}.

Instruction tuning is essential for aligning LLMs with user intentions~\citep{ouyang2022training,cao2023instruction}. Initial efforts in this domain focused on constructing instruction-following datasets through three primary channels: 1) adapting existing NLP benchmarks into instruction formats~\citep{longpre2023flan,wang2022super}; 2) manual curation via crowdsourcing~\citep{DatabricksBlog2023DollyV2, kopf2023openassistant}; and 3) distilling responses from proprietary models~\citep{alpaca, vicuna2023}. While these methods provide high-quality responses, they often fail to capture the complexity of real-world user queries. Therefore, the automated synthesis of complex instructions has emerged as a challenge.

Several strategies have been proposed to address the scarcity of complex instructions. For instance, Evol-Instruct \citep{xu2023wizardlm} employed LLMs to incrementally increase the depth and breadth of instructions through constraints, deepening, and concretizing. Similarly, Conifer \citep{sun2024conifer} utilized a multi-level taxonomy to inject fine-grained constraints into instructions. However, these methods primarily rely on few-shot prompting, which often results in synthetic constraints that diverge from the natural distribution of complex instructions in real applications.

Back-translation, a process of translating text from the target language back into the source language, is mainly used for data augmentation in tasks like machine translation~\citep{sennrich2015improving, hoang2018iterative}.  \citet{li2023self} first applied this to large-scale instruction generation using unlabeled data, with Suri~\citep{pham2024suri} and Kun~\citep{zheng2024kun} extending it to long-form and Chinese instructions, respectively. \citet{nguyen2024better} enhanced this method by adding quality assessment to filter and revise data. Building on this, we further investigated methods to generate high-quality complex instruction datasets using back-translation in an iterative manner.

\subsection{Complex Instruction Alignment}

Alignment methods are essential for enhancing complex instruction-following capabilities. Early approaches primarily relied on Supervised Fine-Tuning (SFT) using data distilled from stronger models \citep{xu2023wizardlm, sun2024conifer}. For instance, WizardLM \citep{xu2023wizardlm} utilized an iterative strategy to construct training data for SFT. I-SHEEP \citep{isheep} introduced an SFT strategy based on an active self-enhancement paradigm where the model aligns itself. However, the inherent limitations of imitation learning often constrain the overall performance of these methods \citep{gudibande2023falsepromiseimitatingproprietary}.

Recent studies also employ Direct Preference Optimization (DPO) on preference pairs, where negative samples are misaligned with the complex instruction \citep{huang2025musc, he2024complex, pham2024suri}. These methods either rely on automatic response detachment \citep{huang2025musc,ren2025step} or rule-based verification \citep{guo2025recast} to construct misaligned responses. However, the efficacy of DPO is often limited because negative samples frequently originate from models with different parameter distributions from the target model, leading to optimization bottlenecks \citep{ivison2024unpacking}. Therefore, other studies also started to employ online reinforcement learning methods to mitigate the distribution drift \citep{ren2026lsriflogicstructuredreinforcementlearning,qin2025incentivizingreasoningadvancedinstructionfollowing}. Nevertheless, the lack of high-quality reward signals remains a major hurdle for online optimization. Developing more efficient strategies for complex instruction alignment remains a critical problem.